\documentclass[../main.tex]{subfiles}
\begin{document}
%==========================================TIÊU ĐỀ TẬP========================================
\begin{table}[h]
    \begin{adjustwidth}{-1cm}{}
        \begin{tabular}{>{\centering\arraybackslash}p{6cm}|>{\raggedright\arraybackslash}p{12.5cm}}
            \multirow{5}{6cm}
            % Cột bên trái: ÉPISODE và số tập
            {\\[-40pt]
                \flaregothic\fontsize{20pt}{20pt}\selectfont \centering
                \color{eptype}ÉPISODE\color{black}\\
                \burbank\fontsize{72pt}{72pt}\selectfont
                \color{epnum}41\color{black}\\[6pt]
                \cafeteria\fontsize{20pt}{20pt}\selectfont\color{epnum}(3-11)\color{black}
                \\[18pt]
                \flaregothic\fontsize{18pt}{18pt}\selectfont
                \color{epattr}ÉPISODE PERDU
                \color{black}
            }
             &
            % Dòng thứ nhất: tên tập tiếng Nhật
            {\fotskip\fontsize{18pt}{18pt}\selectfont 〇〇\ruby{捕}{ほ}\ruby{獲}{かく}\ruby{選}{せん}\ruby{手}{しゅ}\ruby{権}{けん}！
            }                                                                                      \\[6pt]
            % Dòng thứ hai: tên tập tiếng Anh
             & {\cafeteria\fontsize{18pt}{18pt}\selectfont Electric Beach Fever}                                                                                \\[4pt]
            % Dòng thứ ba: tên tập tiếng Việt
             & {\vnmsans\fontsize{18pt}{18pt}\selectfont Cuộc đi săn cua kỳ thú}                                                                                 \\[8pt]
            % Dòng thứ tư: Nhân vật xuất hiện
             & {\caviar{\fontsize{14pt}{14pt}\selectfont Nhân vật: \emp{kuon1} \emp{aqua} \emp{ayame} \emp{okayu}}}
            \\[8pt]
            % Dòng cuối cùng: Thời gian diễn ra của tập (thường sẽ là ngày phát hành)
             & {\continuum\fontsize{16pt}{16pt}\selectfont Ngày phát hành: 16/02/2020}                                                                \\[18pt]
        \end{tabular}
    \end{adjustwidth}
\end{table}
%=========================================TRUYỆN CHÍNH========================================
\color{black}

\txtbx[2]{{\color{vol} \uline{Tóm tắt:}} Một cơn thèm hải sản bột phát của Aqua đã kéo theo Kuon, Ayame và Okayu vào một chuyến đi bắt cua đầy ngẫu hứng ngoài khơi xa. Thay vì mang về những mẻ cua béo ngậy, họ lại liên tục ``cua'' được những thứ không tưởng: từ chiếc càng khổng lồ kẹp bùa may mắn, một quả mận chua biết nói đố mẹo, cho đến sinh vật kỳ dị lai giữa cua và cá sấu với cái kết ``đắng lòng'' sau 100 ngày.}

\begin{center}
  \uline{(Tập này có sử dụng lối chơi chữ đồng âm giữa các ngôn ngữ để tăng tính hài hước.)}
\end{center}

Một buổi chiều trung tuần tháng Hai, cái se lạnh cuối đông vẫn đang còn lưu luyến, bao phủ lấy từng góc nhỏ trong văn phòng \hll. Ba cô nàng Aqua, Okayu và Ayame đang ngồi tụ tập quanh bàn trà, gương mặt ai nấy đều lộ rõ vẻ uể oải dưới bầu không khí trầm lắng. Thế nhưng, sự yên bình ấy chẳng kéo dài được lâu. Bất chợt, Aqua bật dậy như một chiếc lò xo bị nén chặt, đôi mắt hải quân sáng rực lên đầy hưng phấn.

Cô nàng giơ nắm tay lên trời, hét lớn để phá tan sự tĩnh lặng: ``Mọi người ơi! Có ai muốn ăn cua (カニ) không nào?'' 

Kuon lúc này đang mải mê kiểm tra lịch trình dày đặc trên điện thoại cũng phải giật mình ngẩng đầu lên. Cậu xoa xoa cằm gật gù tán thành: ``Cua à? Nghe cũng hấp dẫn đấy, lâu rồi anh cũng chưa được thưởng thức hương vị tươi ngon của biển cả.'' 

\begin{figure}[H]
  \centering
  \begin{subfigure}[t]{0.45\textwidth}
    \centering
    \includegraphics[width=8cm]{../../../../Image/Book3/3-12a1.png}
  \end{subfigure}
  \quad
  \begin{subfigure}[t]{0.45\textwidth}
    \centering
    \includegraphics[width=8cm]{../../../../Image/Book3/3-12a2.png}
  \end{subfigure}
\end{figure}

Okayu đang nhâm nhi túi đồ ăn vặt bỗng dừng lại, liếc nhìn Aqua với ánh mắt đầy hoài nghi. Chẳng nói chẳng rằng, cô nàng thản nhiên kéo chỏm tóc của Ayame, đẩy cô bé ra phía trước mặt mình như một món đồ trưng bày triển lãm rồi hỏi: ``Cua hả (かに)?'' 

Aqua tức tối nhăn mặt, xua tay lia lịa phản đối: ``Đó là Quỷ (オニ) mà! Bà lôi Ojou ra làm hình minh họa cái kiểu gì thế hả!?'' 

Bị đem ra làm trò đùa, Ayame chỉ khẽ nghiêng đầu, khuôn mặt hiện rõ vẻ ngơ ngác đáng yêu. Trên tay cô bé lúc này bỗng xuất hiện một quả mận chua tròn xoe màu đỏ cam, trên thân có ghi rõ chữ ``カキ''. Cô bé chớp chớp mắt, hỏi bằng giọng ngây thơ vô số tội: ``Thế cái này mới là cua (かに) sao ạ?'' 

Nhìn thấy thứ trái cây không chút liên quan, cả Kuon lẫn Aqua đều không hẹn mà gặp, đồng thanh hét lên đến khản cả cổ: ``Cái đó là quả mận chua (カキ) mà! Mấy đứa có chịu phân biệt được giữa hải sản tươi sống với hoa quả chín cây không vậy hả trời!?'' 

Bất chấp sự hỗn loạn, Aqua lập tức xoay người, hướng đôi mắt long lanh ngập tràn tia hy vọng về phía Kuon, mong chờ vị Tổng quản lý hào phóng sẽ rút hầu bao chiêu đãi cả nhóm một bữa tiệc cua hoàng đế thịnh soạn. Thế nhưng, đáp lại sự kỳ vọng đó, Kuon chỉ biết cười trừ rồi nhún vai, thọc tay lộn ngược hai chiếc túi quần trống rỗng của mình ra ngoài: ``Đừng nhìn anh bằng ánh mắt đó. Quỹ văn phòng tháng này đã hoàn toàn cạn kiệt sau mấy vụ đập phá tanh bành của Shion rồi. Giờ mà muốn ăn cua ấy thì\dots'' 

Cậu còn chưa kịp dứt câu trần tình, nàng Hầu gái tóc tím đã giơ cao nắm đấm, reo lên với ý chí kiên định: ``Tự túc là hạnh phúc! Nếu không có tiền mua thì chúng ta tự đi bắt cua đi!'' [cite: 20, 21]

Trước tinh thần rực lửa ấy, một kế hoạch ``thám hiểm đại dương'' đầy ngẫu hứng nhanh chóng được thông qua. Họ mượn tạm một con tàu đánh cá nhỏ thuộc quyền sở hữu của công ty, thu gom đủ loại lưới, cần câu cùng bẫy cua chuyên dụng, rồi hào hứng nhổ neo, rẽ sóng hướng thẳng ra vùng khơi xa bạt ngàn.
\img{3-12b} 

Gió biển thổi lồng lộng qua boong tàu. Ba cô nàng ngồi xếp thành một hàng dài, kiên nhẫn thả những sợi dây cước xuống làn nước xanh thẳm, sâu hoắm bên dưới. Ở phía trong khoang lái, Kuon vừa vững tay giữ vô lăng điều khiển hướng tàu, vừa chăm chú quan sát từng chấm tín hiệu trên màn hình radar dò tìm thủy sản. Đã ba tiếng đồng hồ đằng đẵng trôi qua, vầng mặt trời đỏ rực bắt đầu chìm dần xuống đường chân trời, dát một lớp ánh sáng vàng óng ánh rực rỡ lên mặt biển cô quạnh, nhưng tuyệt nhiên vẫn chẳng có lấy một chiếc càng cua nào thèm xuất hiện.

Ayame mệt mỏi thở dài thườn thượt, ngón tay xoay xoay thân chiếc cần câu một cách vô vọng: ``Chẳng có con cua nào chịu cắn câu cả\dots\ Hay là toàn bộ cua ở vùng biển này kéo nhau đi ngủ hết rồi nhỉ?'' 

Okayu lười biếng nhấp một ngụm nước cam, khuôn mặt đột ngột nhăn nhó lại vì vị chua loét: ``Chua quá đi mất\dots\ Chua đến mức đến lũ cua dưới kia chắc cũng sợ mà chạy mất dép rồi chăng?'' 

Trong khi đó, Aqua đang nằm ườn người trên thành boong tàu đầy uể oải. Cô nàng đột nhiên ngẩng đầu lên nhìn bầu trời xa xăm với gương mặt mang đậm nét triết lý sâu xa: ``Này mọi người, tại sao chúng ta lại gọi việc đi bắt cua là 'cua cua' nhỉ? Nghe nó cứ\dots\ lặp đi lặp lại một cách vô lý thế nào ấy.'' 

Giọng nói của Kuon ngay lập tức vang lên qua hệ thống đàm thoại của con tàu, không thể giấu nổi vẻ buồn cười trước suy nghĩ non nớt ấy: ``Anh cũng chịu đấy. Nhưng biết đâu, có lẽ vì đi bắt cua sẽ khiến cho tâm hồn chúng ta trở nên\dots\ cua vui chăng?'' 

Nàng Hầu gái nghe vậy liền phì cười, lớn tiếng trêu chọc: ``Trời ạ, anh giải thích nghe còn 'chua' hơn cả ly nước cam của Okayu nữa đó!'' 

Đúng lúc bầu không khí chuẩn bị chìm vào sự tẻ nhạt, chiếc cần câu trong tay nàng Hầu gái bỗng dưng giật mạnh một cú kinh hoàng, lực kéo dữ dội đến mức suýt chút nữa đã lôi tuột cô nàng xuống làn nước sâu thẳm. Aqua trợn tròn đôi mắt, vội vàng bấm chặt chân xuống boong, dồn hết sức bình sinh để giữ chặt lấy tay cầm và cuộn dây cước, miệng hét lớn vang vọng cả một góc biển: ``\textbf{Lên rồi! Hàng về rồi! Cua lớn đây rồi mọi người ơi!}'' [cite: 29, 30]

Ayame và Okayu nghe thấy tiếng động liền lập tức vứt bỏ vẻ uể oải ban nãy, lao nhanh về phía thành tàu, đôi mắt lấp lánh đầy mong đợi: ``Cố lên Aqua! Kéo mạnh lên! Phải là một con cua thật to nha!'' [cite: 30, 31]

Từ bên trong khoang lái, Kuon cũng vội ló đầu ra ngoài cửa sổ hét vọng xuống đầy tinh quái: ``\textbf{Cẩn thận đấy A-cua! Giữ chắc tay vào, đừng có để 'của' đi thay người đấy nhé!}'' [cite: 31, 32]

Dù đang trong tình huống gay cấn, cả Ayame lẫn Okayu đều phải bật cười khúc khích trước màn chơi chữ đầy ngẫu hứng và có phần vô tri của ông anh Tổng quản lý, giúp bầu không khí căng thẳng trên boong tàu phần nào được giải tỏa. Trong khi đó, Aqua vẫn đang gồng mình giành giật quyết liệt với sinh vật bí ẩn dưới đáy sâu đang cố ghì chiếc cần câu trũng xuống thành một đường cong hoàn hảo. Cô nàng hét lên một tiếng dài đầy dứtkhát, dùng toàn bộ trọng lượng cơ thể dồn lực kéo mạnh một cú cuối cùng. 

Vút! Một vật thể đen ngòm văng mạnh lên không trung, xoay vài vòng điệu nghệ giữa nền trời hoàng hôn rồi rơi bịch xuống boong gỗ. Kuon bằng phản xạ nhạy bén đã lao ra khỏi khoang lái, nhanh tay dùng một chiếc lưới lớn chụp gọn lấy nó ngay lập tức. Thế nhưng, khi cả bốn người cùng tò mò xúm lại gần để kiểm tra thành quả, tất cả bỗng chốc rơi vào trạng thái đứng hình toàn tập.

Thứ đang nằm giãy giụa trong lưới hoàn toàn không phải là một con cua nguyên vẹn. Đó chỉ là một chiếc càng cua khổng lồ, thô ráp, có kích thước còn lớn hơn cả một lòng bàn tay người lớn, và điều dị hợm nhất là nó đang kẹp chặt lấy một lá bùa hộ mệnh màu đỏ rực đầy ma quái.

\begin{figure}[H]
  \centering
  \includegraphics[width=13.5cm]{../../../../Image/Book3/3-12c.png}
  \caption{Trong phiên bản gốc, con cua đó đang kẹp vào một cái ủng hiệu Kani x Kani (do nó vần với từ ``kani'' trong tiếng Nhật).} 
\end{figure}

Ayame ngơ ngác nghiêng đầu, nhìn chằm chằm vào chiếc càng đơn độc: ``Cái này\dots\ cũng được tính là cua hả mọi người\dots?'' 

Okayu tò mò cúi xuống nhặt vật thể kỳ dị ấy lên, xoay qua xoay lại quan sát dưới góc độ của một chuyên gia sinh vật học đầy chuyên nghiệp rồi đưa ra lời chốt hạ xanh rờn: ``Theo phân tích của tôi, thứ này chứa khoảng 30\% gen của loài cua, còn 70\% còn lại hoàn toàn là cấu trúc của một linh vật tâm linh huyền bí.'' 

Kuon bất lực nhếch môi nở một nụ cười đau khổ: ``Anh thì chỉ thấy đúng một cái càng bám đầy rong rêu với cái bùa chú quái đản kia thôi. Còn bản thể con cua tội nghiệp kia chắc giờ này đang phải lội bộ về bờ để báo cảnh sát vì bỗng dưng bị giật mất cánh tay rồi đấy.'' 

Aqua chán nản lắc đầu đầy thất vọng trước món bùa chú xui xẻo này. Cô nàng dứt khoát giật lấy chiếc càng quái dị, thẳng tay ném nó trở lại lòng đại dương bao la: ``Không tính! Cái này không ăn được! Tiếp tục câu con khác thôi!'' [cite: 39, 40]

Nhìn chiếc càng chìm dần xuống dòng nước xiết, Okayu liền cầm chiếc càng cùng cái bùa quăng trở lại biển, mặc cho tiếng thở dài tiếc nuối của nàng Oni: ``Mất bao nhiêu công sức kéo lên thế mà lại vứt đi uổng vậy sao\dots'' 

Như để đáp lại sự kiên trì của họ, chỉ vài phút ngắn ngủi sau đó, chiếc cần câu trong tay Aqua lại một lần nữa rung lắc dữ dội. Nàng Hầu gái hăm hở dùng hết sức kéo ngược dây cước lên trong sự kỳ vọng lên tới tột đỉnh: ``\textbf{Được rồi! Lần này thần may mắn phải mỉm cười với chúng ta, chắc chắn nó phải là một con cua xịn!}'' 

Nghe tiếng gọi reo, cả ba người còn lại lập tức bỏ dở việc đang làm, vây quanh Aqua và nín thở chờ đợi giây phút sinh vật biển lộ diện. Thế nhưng, đến khi vật thể lạ hoàn toàn được vớt lên khỏi mặt nước, toàn bộ boong tàu lại một lần nữa rơi vào bầu không khí im phăng phắc đầy ngỡàng. 

Đó là một quả mận chua tròn trịa, được thắt một chiếc nơ ruy băng màu đỏ vô cùng xinh xắn ở cuống. Điều kinh dị và quái đản đến mức rợn tóc gáy là quả mận ấy bỗng nhiên rung rinh, mở ra một cái khe nhỏ như chiếc miệng rồi cất tiếng nói bằng một chất giọng eo éo, chói tai: ``Hỏi đố cả đám nè: Tại sao tổ chức \hll\ lại được gọi là \hll\ vậy ta?'' 

\begin{figure}[H]
  \centering
  \includegraphics[width=13.5cm]{../../../../Image/Book3/3-12d.png}
  \caption{Trong phiên bản gốc, thứ mà Aqua câu được là một quả hồng có ghi chữ カキ ở trên trán. Câu hỏi mà nó hỏi mọi người cũng tương tự trong tiếng Việt.} [cite: 44, 45]
\end{figure}

Đối diện với câu hỏi hóc búa từ một sinh vật thực vật kỳ dị, cả bốn người đứng nhìn nhau trân trân, không ai biết nên phản ứng ra sao. Gương mặt Aqua lúc này đã tối sầm lại vì tức giận, cảm giác như lòng tự trọng của một nàng hầu gái tối cao vừa bị một quả trái cây nhỏ bé vô danh xúc phạm và ``kháy đểu'' một cách trắng trợn. Chẳng thèm suy nghĩ lấy một giây, cô nàng lập tức vung tay, thẳng thừng ném mạnh quả mận nói nhiều ấy bay một đường parabol tuyệt đẹp xuống dưới biển sâu: ``Đi mà hỏi lũ cua dưới kia ấy! Đồ quả mận chua ngoa vô duyên thối nát kia!'' [cite: 47, 48]

Đến lần thứ ba này, Aqua quyết định chơi tất tay với vận may của mình. Cô nàng bậm môi, gồng toàn bộ cơ bắp ở tay để kéo chiếc cần câu với một lực kinh hoàng phát ra tiếng gió vút vút. Và rồi, khi mặt nước rẽ ra, thứ vừa lộ diện dưới ánh hoàng hôn tà ác đã khiến cho tất cả mọi người có mặt trên tàu đều phải rùng mình, nổi hết cả da gà. Đó là một sinh vật lai tạp vô cùng dị hợm, sở hữu phần đầu và hai chiếc càng to đùng của một con cua khổng lồ, nhưng phần thân sau lại kéo dài thành một cái đuôi thượt cùng lớp da gai góc, sần sùi mang đậm đặc trưng của một con cá sấu đầm lầy.
\img{3-12e} 

Ayame run rẩy xáp lại gần, khẽ lẩm bẩm đầy hoang mang: ``Cái quái thai gì thế này\dots\ Đây là loài cua sấu hay là cá cua vậy mọi người?'' 

Okayu đứng khoanh tay, chăm chú nhìn chằm chằm vào sinh vật tội nghiệp đang nằm bò trên sàn tàu bằng ánh mắt vô hồn, rồi buông ra một lời tuyên bố lạnh lùng đến đáng sợ: ``Theo kinh nghiệm tâm linh nhiều năm của tôi, chính xác là 100 ngày sau con này kiểu gì cũng chết chắc\footnote{Liên hệ tới \textit{The Alligator who will Die in 100 Days} (100 ngày sau, con cá sấu sẽ chết).}.'' 

Kuon cũng khoanh tay đứng bên cạnh, nghiêng đầu đưa ra lời đánh giá mang tính chuyên môn cao: ``Anh thì bi quan hơn, chắc chỉ cho nó sống thọ được khoảng 60 ngày thôi. Nhìn sắc mặt nó có vẻ đang bị stress tâm lý nặng nề sau khi bị chiếc cần câu của chúng ta kéo thốc lên từ độ sâu trăm mét rồi.'' 

Trái ngược với vẻ u ám của hai người kia, Ayame vẫn giữ vững tinh thần lạc quan vốn có. Cô bé chỉ tay vào chiếc mai cứng cáp của sinh vật dị dạng: ``Nhìn kỹ mà xem, rõ ràng có tới 70\% cơ thể của nó mang đặc tính của cua mà! Cứ tạm tính nó là một con cua đi thôi!'' [cite: 51, 52]

``Nếu tính như vậy thì chúng ta còn khoảng 30 ngày để tìm cách chế biến và ăn thịt nó,'' Okayu dửng dưng tiếp lời.

Kuon nhún vai cười trừ: ``Anh thì đoán chỉ cần 10 ngày là nó lên thớt rồi. Akutan, em là trưởng đoàn, em quyết định số phận của nó thế nào đây?'' 

Aqua nhìn sinh vật lai tạp đang nằm thoi thóp, thở không ra hơi trên boong tàu. Cô nàng bỗng nở một nụ cười đầy tự tin, đập tay một cái bốp: ``OKua! Cứ xách cổ mang nó về văn phòng đi, dù sao nhìn cái thớ thịt sần sùi kia của nó\dots\ trông cũng có vẻ khá là ngon miệng đấy chứ!'' [cite: 53, 54]

Như hiểu được số phận thảm khốc đang chờ đón mình ở tương lai, con cua lai cá sấu ban nãy còn đang hấp hối bỗng dưng bật dậy như một vị thần. Nó vung vẩy hai chiếc càng, hét lớn bằng giọng phẫn uất tột cùng: ``Mấy người kia có cái tính người không hả, định không thả tôi về với gia đình luôn đúng không!?  Tôi thề tôi chỉ là một con cua tội nghiệp vô tình đi lạc vào cái dòng hải lưu quái quỷ này thôi mà!  Đừng có tự tiện mang tôi về để làm món lẩu cua chứ!!!'' 

Tiếng thét thất thanh rền vang đầy căm phẫn của sinh vật kỳ lạ khiến Aqua giật bắn mình, nhảy dựng ra sau vì hoảng hốt. Trong khi đó, ba người còn lại thì không thể nhịn nổi nữa, liền ôm bụng cười ngặt nghẽo trước màn đối thoại vô tiền khoáng hậu này. Nhưng dù sinh vật đó có van xin thảm thiết đến đâu, cả bốn người rốt cuộc vẫn quyết định tạm thời giữ chú ``cua lai'' tội nghiệp này lại để đem về tiến hành nghiên cứu sinh học thêm. Cuối cùng, cuộc hành trình đi săn cua đầy bão táp đã khép lại theo một kịch bản không thể nào mang đậm phong cách ``hologra'' hơn thế.

%==========================================TRUYỆN PHỤ=========================================
\omk

\begin{center}
  \textit{\uline{Ngày hôm sau.\\Tại một góc sân trường Đại học.}}\\
  \textit{*Các nhân vật nói bằng tiếng Etsunan.}
\end{center} 

Vào giờ nghỉ giải lao ngày hôm sau, dưới bóng râm mát mẻ của một góc sân trường Đại học đầy nắng, Kuon đang tụ tập quanh đám bạn thân nối khố trong lớp để hào hứng kể lại chiến tích đi bắt cua vô tiền khoáng hậu của mình[cite: 59, 60]. Cậu chàng vừa thuật lại vừa nhiệt tình khua tay múa chân một cách sinh động để mô tả lại hình thù quái dị của những sinh vật mà mình đã chạm trán giữa biển khơi sâu thẳm[cite: 60, 61].

Nghe câu chuyện có phần hoang đường, Shunyo tò mò rướn người tới trước hỏi: ``Thế rốt cuộc là chuyến đi hoành tráng đó ông có bắt được mẻ hải sản nào ra ngô ra khoai không, hay lại xách giỏ không đi về?'' 

Kuon phì cười lớn, vỗ đùi đánh bôm bốp đầy đắc ý: ``Cua thì phải nói là chất thành núi, mang về ăn ba đời không hết! Nhưng mà cái giống cua ở vùng biển đó nó lạ lùng và kinh dị lắm các ông ạ!  Đầu tiên là cái càng cua kẹp bùa chú tâm linh, xong rồi đến một quả mận chua biết mở mồm nói tiếng người đố mẹo, rồi đỉnh điểm nhất phải kể đến con quái vật đầu cua đuôi cá sấu cơ!''  Cậu con trai say sưa liệt kê ra một danh sách dài dằng dặc những chủng loại sinh vật đột biến mà bản thân cùng ba cô nàng idol đã câu được trong suốt ngày hôm đó, khóe miệng vẫn không thể nào ngừng cười toe toét, ``Xong rồi còn có cả cua chân vịt nữa chứ, mà chưa dừng lại ở đó đâu, còn cả một đống thứ cua dị dạng phát quang khác nữa nhìn mà hãi hùng luôn!'' 

``Cái quái gì vậy thật sao\dots!? Ông vừa dắt đám con gái người ta đi câu ở vùng biển rác thải hạt nhân hay khu phóng xạ nào đấy à?'' Hokuro thảng thốt kêu lên, mắt trợn trừng kinh hãi[cite: 63, 64].

``Tôi cũng chịu chết thôi, chỉ biết là cái vùng nước chỗ đó rêu xanh phủ kín rì, lũ cua thì bò lổm ngổm dày đặc đè lên nhau cứ như thể đang tụ tập biểu tình đòi quyền sống vậy,'' Kuon cười lớn thành tiếng rồi khom lưng xách ra bốn chiếc túi canvas to đùng, nặng trịch đặt lên mặt bàn đá. ``Vì bắt được nhiều quá ăn không xuể nên tôi mang hẳn sang đây để biếu cho cả lớp một ít làm quà này!'' 

Nghe thấy thế, đám con trai trong lớp lập tức xúm đông xúm đỏ lại. Ai nấy đều mắt tròn mắt dẹt, há hốc mồm kinh ngạc nhìn chằm chằm vào đống cua dị dạng đang bò lổm ngổm bên trong: có con thì sở hữu chiếc càng khổng lồ to và sắc bén như một cái kìm của robot công nghiệp, có con thì phần vỏ xù xì bám đầy những gai nhọn hoắt trông cực kỳ nguy hiểm, lại có con rực rỡ đủ loại sắc màu neon phản quang cứ như thể vừa mới bước ra từ một bộ phim hoạt hình viễn tưởng vậy.

``Cái thứ quái đản này mà ông cũng mặt dày bảo là cua được á? Nhìn nó có khác gì một sinh vật ngoại lai trộn lẫn với\dots\ cái quái thai gì ấy nhỉ?'' Fukuyasu vừa lia máy điện thoại chụp ảnh lia lịa vừa lầm bầm đầy hoài nghi.

``Cua lai quái vật đột biến gen thì có chứ cua bình thường nào thế này!'' Một cậu bạn khác đứng bên cạnh nhanh nhảu chêm vào một câu.

Kuon đưa tay lên gãi gãi đầu cười ngượng ngùng trước phản ứng dữ dội của đám bạn: ``Thì dù hình thù nó có hơi phá cách một chút nhưng xét về mặt sinh học thì nó vẫn là cua mà!  Tối nay ông nào rảnh rỗi thì cứ qua thẳng nhà tôi, tôi đích thân xuống bếp làm một nồi lẩu cua thập cẩm đại bổ để đãi cả lớp một bữa ra trò. Đảm bảo hương vị thơm ngon đến mức quên sạch mọi sầu đời luôn!'' 

Nghe tới việc được ăn lẩu miễn phí do chính tay vị Tổng quản lý đảm nhận nấu nướng, cả đám con trai liền hò reo vang dội, vỗ tay hưởng ứng đầy nhiệt liệt, ai nấy đều vô cùng hào hứng chờ đợi đến buổi tối hẹn hò.
\ct 

Trong lúc cả hội đang mải mê bàn tán rôm rả về danh sách thực đơn đi kèm cho món lẩu tối nay, Tokue bất chợt huých vai Kuon, nháy mắt một cái đầy ẩn ý trêu chọc: ``Mà này ông bạn, tôi nghe phong thanh có tin đồn là hôm qua ông không chỉ đơn thuần đi bắt cua biển đâu nha, mà còn tranh thủ 'cua' luôn được mấy em gái xinh tươi mơn mởn đi cùng chuyến tàu nữa đúng không hử?'' 

``Cua cái đầu ông ấy mà cua!'' Kuon tức mình cốc nhẹ một cái rõ đau vào đầu cậu bạn thân để cảnh cáo. ``Tôi là đi khảo sát thực tế và hỗ trợ công việc cùng với ba người đồng nghiệp trong công ty, chứ có phải là đi hẹn hò yêu đương nhăng nhít gì đâu mà ông cứ vẽ chuyện.'' 

``Ơ hay, thế ba người đồng nghiệp đó cụ thể là những ai mà bấy lâu nay ông cứ giấu kỹ như mèo giấu cứt thế?'' Miharu chen ngang với giọng điệu đầy vẻ hóng hớt, tai dựng đứng lên chờ đợi.

``Thì gồm có một cô nàng Quỷ Oni lai thần linh, một cô Mèo lém lỉnh thích ăn vụng, và một nàng Hầu gái tóc tím buộc tóc hai bím xinh xắn,'' Kuon thản nhiên mở miệng liệt kê một tràng ngắn gọn.

Nghe đến những đặc điểm nhận dạng vô cùng quen thuộc và độc nhất vô nhị đó, cả đám bạn lập trình xung quanh lập tức im bặt trong một giây, rồi ngay sau đó đồng loạt rộ lên tiếng ``Ồ~'' kéo dài đầy ranh mãnh khi nhận ra ngay lập tức những thần tượng đình đám mà Kuon vừa thản nhiên nhắc tới chính là dàn nhân vật chủ chốt thuộc biên chế của tập đoàn giải trí \hll.

``Này\dots\ Ông bạn già ơi, ông có bằng chứng xác thực nào cho lời mình vừa nói không đấy? Đừng có mà thấy tụi này hiền rồi đứng đây bốc phét, nổ banh xác nha nha!'' Tokue khoanh tay tỏ vẻ vô cùng nghi ngờ.

Kuon khẽ nhếch môi nở một nụ cười đắc thắng đầy tự mãn. Cậu từ tốn lôi từ trong túi áo khoác ra hai tấm ảnh vừa mới in xong xuôi: ``Tang chứng vật chứng rõ ràng rành rành đây, xem đi rồi hẵng phán nhé mấy ông tướng!'' 

Tấm ảnh đầu tiên chụp lại cận cảnh hai cô nàng Ayame và Okayu đang cặm cụi, ghé sát đầu vào nhau để chơi trò\dots\ bầu cua cá tôm ăn tiền trên một cái bảng giấy cũ nát, rách rưới đặt ngay giữa boong tàu với những biểu cảm trên gương mặt cực kỳ căng thẳng và nghiêm túc như thể đang tham gia trận chung kết thế giới. 

Còn tấm ảnh thứ hai thậm chí còn chấn động hơn, ghi lại cảnh tượng Aqua đang hùng hổ dùng một sợi dây thừng cỡ lớn để trói chặt một ông\dots\ cua-rơ điền kinh thực thụ bằng xương bằng thịt đang mặc nguyên bộ đồ thể thao bó sát, không ngừng giãy giụa đầy đau khổ trên sàn tàu.

Cả đám bạn đứng hình toàn tập, há hốc mồm mất vài giây trước mức độ phi lý và điên rồ đến tận cùng của hai bức ảnh chụp.

Tokue lắp bắp, ngón tay run rẩy chỉ vào tấm ảnh: ``Thế\dots\ Thế cái quái gì đây\dots\ Làm cách nào mà giữa đại dương bao la các ông lại có thể dùng cần câu để vớt được cả một vận động viên cua-rơ lên trên tàu vậy hả trời!?'' 

Kuon khẽ thở dài, bước tới vỗ vỗ vai cậu bạn thân của mình, buông một câu chốt hạ mang đầy sự trải đời cay đắng: ``Vì cái nơi đó chính là văn phòng \hll\ mà ông bạn. Ở cái chốn kỳ lạ đấy thì logic với định luật vật lý chỉ là những khái niệm vô cùng xa xỉ mà thôi!'' 

Nhận được câu trả lời không thể thuyết phục hơn, cả đám con trai lại được một trận cười vỡ bụng vang dội khắp cả góc sân trường, khép lại câu chuyện về chuyến hành trình ``đi câu cua'' đầy huyền thoại mà chắc chắn sẽ còn được lưu truyền mãi trong biên niên sử của lớp học nhà Kuon.
\ct 

Sau bữa tiệc lẩu cua ``thập cẩm'' diễn ra vô cùng đại thành công vào tối hôm đó cùng đám chiến hữu, Kuon quay trở lại văn phòng làm việc vào hôm sau với một tâm trạng vô cùng sảng khoái và tràn đầy năng lượng. Trong giờ nghỉ giải lao, cậu thong thả nhấp ngụm trà rồi hào hứng kể lại toàn bộ phản ứng đầy hài hước của đám bạn học cho ba cô nàng ``tội đồ'' đầu têu chuyến đi hôm trước cùng nghe[cite: 82, 83].

``Mấy đứa không thể tưởng tượng nổi đâu, đám bạn trong lớp anh ăn mấy con cua biến dị đó mà đứa nào đứa nấy cứ tấm tắc khen ngon lấy khen để, thậm chí còn gạ gẫm đòi lần sau anh phải dắt tụi nó đi bắt tiếp nữa kìa!'' 

Okayu nghe vậy liền tỏ ra vô cùng ngạc nhiên, hai tai khẽ động đậy: ``Thật vậy sao ạ?  Mấy cái thứ hình thù kỳ dị quái đản đó mà con người cũng có thể tiêu hóa và khen ngon được á? Ban đầu em cứ đinh ninh mang đống đó về thì chỉ có nước vứt vào tủ kính để làm kiểng trưng bày thôi chứ.'' 

``Nhưng mà\dots\ Chắc là sau khi ăn xong xuôi thì mọi người vẫn hoàn toàn ổn thỏa chứ anh Kuon?'' Ayame lo lắng chớp mắt hỏi, trong lòng có chút bất an về vấn đề an toàn thực phẩm.

Kuon tự tin xua xua tay, trấn an cô bé: ``Yên tâm đi Ojou, tay nghề nấu lẩu của anh đây thì lúc nào cũng là đỉnh của chóp rồi mà!  Sáng nay tụi nó đứa nào đứa nấy vẫn khỏe re như vạm vỡ, đi học điểm danh đầy đủ không thiếu một mống nào cả đâu.'' 

Aqua nghe thấy thế thì đôi mắt hải quân lại một lần nữa sáng rực lên đầy hào hứng: ``Tuyệt vời quá đi mất! Vậy thì lần sau nếu có dịp đi bắt cua ngoài khơi xa thì nhất định, nhất định phải rủ em đi cùng nữa nhé!  Em xin thề và hứa danh dự là lần tới sẽ câu được những con cua hoàng đế thực sự chân chính\dots\ chứ quyết không mang về mấy thứ bùa chú quái dị hay quả mận chua biết nói kia nữa đâu!'' 

Okayu bên cạnh liền bật cười khẩy, buông lời chọc ghẹo quen thuộc: ``Cái bài hứa suông này của bà là lần thứ n trong tháng rồi đấy nhé Akutan~  Để rồi xem, lần sau mà bà có lỡ tay câu được một con cá mập lai cua bạo chúa lên boong thì tự mình đứng ra mà giải quyết hậu quả đấy nhé!'' 

Aqua nghe vậy liền phồng má, trợn mắt giận dỗi quay mặt đi chỗ khác, khiến cho bầu không khí trong cả căn văn phòng lại một lần nữa rộn rã tiếng cười đùa vô tư. Dù cho chuyến hành trình thám hiểm đại dương có trôi qua theo một cách kỳ quặc, điên rồ và nằm ngoài mọi quy luật logic đến đâu, thì đối với mỗi người họ, đó vẫn luôn là những mảnh ký ức lấp lánh, ngập tràn niềm vui tươi đẹp của tuổi trẻ.
\end{document}